
au degré \(0\).

Le module des différentielles \(\underline{\Omega}^1_{\underline{B}/\underline{O}_P}\) est isomorphe à \(\underline{B}(-1)\otimes_{\underline{O}_P}\underline{O}_P(1)\) car on peut définir une \(\underline{\underline{S}}(\underline{E})\)-dérivation \(d\) :
\[
\underline{\underline{S}}(\underline{E})[T,T^{-1}]
\longrightarrow
\underline{\underline{S}}(\underline{E})[T,T^{-1}]\otimes_{\underline{\underline{S}}(\underline{E})}\underline{\underline{S}}(\underline{E})\cdot dT
\]
de façon à satisfaire la propriété universelle qui définit \(\underline{\Omega}^1_{\underline{B}/\underline{O}_P}\) en posant \(d(aT^n)=naT^{n-1}\otimes dT\) ; de plus, \(d\) est compatible aux graduations si on attribue le degré \(1\) au facteur \(\underline{\underline{S}}(\underline{E})\cdot dT\).

Soit \(V\) le \(Y\)-fibré vectoriel défini par \(\underline{E}\). Alors \(\underline{\Omega}^1_{Z/Y}\) est la restriction à \(Z\) de \(\underline{\Omega}^1_{V/Y}\). Or \(V\) est le \(Y\)-schéma affine défini par \(\underline{\underline{S}}(\underline{E})\), donc \(\underline{\Omega}^1_{V/Y}\) est défini par \(\underline{\underline{S}}(\underline{E})\otimes_{\underline{O}_Y}\underline{E}\) (car \(\underline{E}\) est localement libre de type fini), donc \(\underline{\Omega}^1_{V/Y}\) est égal à \(g^*(\underline{E})\), et par suite, \(\underline{\Omega}^1_{\underline{B}/Y}=f^*(\underline{E})\otimes_{\underline{O}_P}\underline{B}\). Pour voir sa graduation, on regarde localement, et on voit aussitôt que \(f^*(\underline{E})\) doit être considéré comme étant de degré \(1\).

En prenant les termes de degré \(0\) dans la suite exacte (5.2.1), on trouve la suite exacte annoncée.

\textbf{Remarque 5.2.1.} En utilisant les résultats généraux sur l'étude différentielle des foncteurs de Hilbert ([2], 5), on obtient directement un isomorphisme \(\underline{\Omega}^1_{P/Y}\xrightarrow{\sim}\underline{F}(-1)\) d'après la caractérisation fonctorielle donnée de \(g^*(\underline{\Omega}^1_{P/Y})\) pour tout \(Y\)-morphisme \(g:Y'\to P\).

\textbf{5.3.} D'après 5.2, on a :
\[
 T_f=\cl(\underline{\Omega}^1_{P/Y})=g^*(E)\cdot\xi^{-1}-1.
\]
Par conséquent :
\[
 \widetilde c(T_f)=\widetilde c(E)\cdot\widetilde c(\xi^{-1})-\widetilde c(1).
\]

